%(BEGIN_QUESTION)
% Copyright 2006, Tony R. Kuphaldt, released under the Creative Commons Attribution License (v 1.0)
% This means you may do almost anything with this work of mine, so long as you give me proper credit

Hastigheten for varmeoverføring ved \textit{stråling} fra et varmt legeme kan uttrykkes med Stefan-Boltzmann-likningen:

$${dQ \over dt} = e \sigma A T^4$$

\noindent
Der,

$dQ \over dt$ = varmeflythastighet

$e$ = emissivitetsfaktor

$\sigma$ = Stefan-Boltzmann-konstanten (5.67 $\times$ $10^{-8}$ W / m$^{2}$ $\cdot$ K$^{4}$)

$A$ = arealet av den strålende overflaten

$T$ = absolutt temperatur

\vskip 10pt

Basert på enheten som er oppgitt for Stefan-Boltzmann-konstanten, bestem riktige måleenheter for varmeflyt, emissivitet, areal og temperatur.

\underbar{file i00343}
%(END_QUESTION)





%(BEGIN_ANSWER)

$dQ \over dt$ = varmeflythastighet (watt)

$e$ = emissivitetsfaktor ({\it enhetsløs})

$A$ = areal av strålende overflate (kvadratmeter)

$T$ = absolutt temperatur (Kelvin)

\vskip 10pt

Utfordringsoppgave: en mer komplett utgave av Stefan-Boltzmann-likningen tar med temperaturen i omgivelsene til det varme objektet:

$${dQ \over dt} = e \sigma A (T_1^4 - T_2^4)$$

\noindent
Der,

$T_1$ = objektets temperatur

$T_2$ = omgivelsestemperatur

\vskip 10pt

Forklar hvorfor dette andre $T$-leddet er nødvendig for at likningen skal gi fysisk mening.

%(END_ANSWER)





%(BEGIN_NOTES)

Svar på utfordringsoppgaven: Hvis vi tolker den første (enklere) likningen bokstavelig, måtte vi enten konkludert med at alle varme objekter er evige energikilder (som stadig stråler ut energi uten å kjøles ned), eller at alle varme objekter vil kjøle seg selv til absolutt null i vakuum (der det ikke finnes ledning eller konveksjon). Med andre ord sier den enkle likningen at objekter bare kan avgi strålingsenergi, ikke motta noe.

%INDEX% Physics, heat and temperature: Stefan-Boltzmann equation

%(END_NOTES)
